% !TEX program = xelatex
% Шаблон научной статьи НИТУ «МИСиС»
% Заполните блоки ниже и скомпилируйте: xelatex article.tex

\documentclass{misisarticle}

% --- Классификационные коды (оставьте пустым {}, если не требуются) ---
\grnti{28.23.15}
\asjc{1707}
\oecd{1.2.01}

% --- Название ---
\title{Cемантической сегментации дефектов на SEM-снимках на основе алгоритмов глубого обучения}

% --- Авторы: \misisauthor{ФИО}{аффилиация}{город, страна}{email} ---
\misisauthor{Н.~А.~Гущин}{кафедра инженерной кибернетики НИТУ «МИСиС»}{Москва, Россия}{m2311537@edu.misis.ru}
\misisauthor{И.~Ю.~Федоров}{кафедра инженерной кибернетики НИТУ «МИСиС»}{Москва, Россия}{}

% --- Аннотация и ключевые слова ---
\abstract{%
	Задача заключается в исследовании подходов для семантической сегментации дефектов на SEM-снимках,применяемых для инспекции полупроводниковых структур. Анализируется влияние архитектуры нейронной сетей и функции потерь на качество выделения дефектов. Рассматриваются архитектуры U-Net,U-Net++ и SegFormer, а также функции потерь Dice loss, focal loss и их комбинация.Дополнительно проводится дообучение HRNet-W18 с обучением головы бинарной сегментации и исследуется вариант с дополнительной головой классификации наличия дефектов. Для обучения и оценки используются открытый набор данных Carinthia-S и изображения из открытых источников, приведённые к бинарной классификации. Сравнение моделей проводится по метрикам mIoU, Dice, Recall и Precision, а также по количеству параметров, FLOPs и времени инференса. Для оценки устойчивости результатов рассчитываются средние значения и стандартные отклонения метрик.
}

\keywords{%
	Система инспекции, кремниевые пластины, SEM-снимки,дефекты полупроводниковых структур,семантическая сегментация
	
}

\begin{document}
	
	\maketitle
	
	% ============================================================
	\section{Введение}
	
	Контроль дефектов является важным этапом производства полупроводниковой продукции,так как нарушения в структуре способны снизить качество и надёжность изделий. Сканирующие электронные микроскопы дают возможность получать изображения поверхности высокого разрешения для нахождения изъянов. Малые размеры дефектов,многообразие их форм и сходство с фоном усложняют анализ снимков.Семантическая сегментация позволяет получать координаты изъянов,что является одной из проблем для автоматической инспекции полупроводниковых структур.Автоматизация систем инспекции позволяет сократить время для обработки изображений и уменьшить человеческий труд.
	
	% ============================================================
	\section{Анализ работ}
	
	В работе Ronnenberger и соавторов~\cite{ronneberger2015} предложена архитектура U-Net,объединяющая декодер и энкодер через пропускающие связи. Передача признаков между разными уровнями помогает улучшить точность распознавания обектов.В данной статье U-Net используется в качестве базовой модели.
	
	В работе Zhou и соавторов~\cite{zhou2018} предложена архитектура U-Net++ с вложенными плотными связями между декодером и энкодером.Такая структура позволяет лучше распознавать объекты,но при этом усложняются вычисления.
	
	В работе Xie и соавторов~\cite{xie2021} предложена архитектура SegFormer,включающая в себя трансформерный энкодер и легкий декодер на основе многослойного перцептрона.
	
	В работе Wang и соавторов~\cite{wang2020} предложена архитектура HRNet, использующая параллельные ветви, сохраняющая высокое разрешение на протяжении всей обработки изображений,а также одновременно обрабатывающая карты признаков. 
	
	В работе Lee и соавторов~\cite{lee2024} предлагают сегментацию дефектов на SEM-снимках с помощью рекурсивного обучения и псевдомасок.Данный вариант обучения снижает затраты на разметку,однако исследуется лишь на 1 типе дефектов.
	
	В работе Hasić и соавторов~\cite{hasic2025} предлагают единый набор SEM-снимков с бинарными масками,который прошел экспертную проверку.Данная подборка изображений
	
	\section{Постановка задачи}
	
	На вход модели подается SEM-снимок полупроводниковой структуры.Требуется получить бинарную маску,в которой каждому пикселю присваивается один из двух классов:дефект или фон.Для модели с дополнительной головой бинарной классификации дополнительно требуется предсакзать наличие дефекта.
	Цель работы - сравнить нейросетевые методы бинарной сегментации дефектов на SEM-снимках и оценить их переносимость на другие данные.
	Для достижения поставленной цели необходимо решить следующие задачи:
	\begin{enumerate}
		\item Подготовить данные и разделить на обучающую,валидационную и тестовую выборку;
		\item Дообучить U-Net,U-Net++,SegFormer и HRNet-W18,а также HRNet-W18 с дополнительной головой бинарной классификации;
		\item Сравнить влияние Dice Loss,Focal Loss и их комбинации на качество сегментации
		\item Оценить качество по метрикам mIou,Dice,Recall и Precision,вычислительные затраты - число параметров и FLOPs и устойчивость результатов по средним значениям и стандартным отклонениям
		\item Для каждой архитектуры выбрать конфигурацию обучения с наибольшим значением Dice на валидационной выборке и оценить её на тестовой выборке и внешнем наборе данных без дополнительного обучения
	\end{enumerate}
	
	\section{Описание методов}
	
	\subsection{Общая схема сегментации}
	
	Перед подачей в модель снимки приводятся к размеру 512$\times$512 пикселей и нормализуются по статистике из ImageNet. Одноканальные SEM-снимки преобразуются в трёхканальные формат RGB путём дублирования значений яркости во всех 3 каналах для согласования с входом предобученных моделей.
	
	\subsection{Архитектуры моделей}
	
	Для сравнения методой выбраны U-Net, U-Net++, SegFormer и HRNet-W18, представляющие разные способы извлечения и объединения информации из снимка.
	
	U-Net~\cite{ronneberger2015} состоит из энкодера и декодера. Энкодер уменьшает пространственное разрешение и увеличивает карту признаков изображения, а декодер восстанавливает исходные размеры для получения попиксельной вероятности.Пропускающие связи передают признаки между уровнями, позволяя использовать пространственные детали при восстановлении маски. В работе применяется U-Net с энкодером Resnet34.
	
	U-Net++~\cite{zhou2018} улучшает архитектуру U-Net за счёт вложенных плотных связей между энкодером и декодером. Дополнительные свёрточные блоки преобразуют признаки перед их объединением. Используется энкодер ResNet34.
	
	SegFormer~\cite{xie2021} объединяет трансформерный энкодер и декодер на основе многослойного перцептрона. Энкодер формирует признаки нескольких масштабов, которые объединяются для получения карты сегментации. В экспериментах используется энкодер MiT-B0.
	
	HRNet-W18~\cite{wang2020} сохраняет разрешения на протяжении обработки изображения. Архитектура содержит параллельные ветви разных разрешений, между которыми осуществляется обмен информацией. Для бинарной сегментации к извлечённым признакам подключается голова,состоящая из Conv(270->128)->BatchNorm->ReLU->Conv(128->1), формирующая карту вероятностей дефекта.
	
	Для инициализации моделей U-Net и U-Net++ используются веса ResNet34,предобученного на ImageNet.Для Segformer веса MiT-B0,а для HRNet предобученные на ImageNet веса. В процессе обучения обновляются все параметры.
	
	\subsection{Дополнительная голова классификации}
	
	Дополнительно исследуется вариант HRNet-W18 с головой бинарной классификации наличия дефекта. Она получает на вход объединенный вектор признаков,полученный по каждой ветви. На выходе формируется вероятность наличия хотя бы одного дефекта на изображении.
	
	Метка равна единице, если маска содержит пиксели дефекта, и нулю в противном случае. Таким образом, сегментационная голова определяет расположение дефекта, а классификационная оценивает его наличие на уровне изображения.
	
	Совместное дообучение выполняется с функцией потерь:
	\[
	L = L_{\mathrm{seg}} + 0.2* L_{\mathrm{cls}},
	\]
	где $L_{\mathrm{seg}}$ - функция потерь сегментации, $L_{\mathrm{cls}}$ - бинарная кросс энтропия.
	
	\subsection{Функции потерь сегментации}
	
	Сравниваются Dice Loss, Focal Loss и их комбинация. Пусть $p_i$~--- предсказанная вероятность дефекта для пикселя $i$, $y_i \in \{0,1\}$~--- соответствующая метка, а $N$~--- число пикселей.
	
	Dice Loss оценивает перекрытие предсказания и истинной маски:
	\[
	L_{\mathrm{Dice}}
	=1-\frac{2\sum_{i=1}^{N}p_i y_i+\varepsilon}
	{\sum_{i=1}^{N}p_i+\sum_{i=1}^{N}y_i+\varepsilon},
	\]
	где $\varepsilon$ - константа, обеспечивающая устойчивость вычислений.
	
	Focal Loss уменьшает вклад распознанных пикселей и усиливает внимание к сложным примерам:
	\[
	L_{\mathrm{Focal}}
	=-\frac{1}{N}\sum_{i=1}^{N}\left[
	y_i(1-p_i)^2 \log p_i+
	(1-y_i)p_i^2\log(1-p_i)\right],
	\].
	
	Комбинированная функция имеет вид
	\[
	L_{\mathrm{seg}}
	=0.5*L_{\mathrm{Dice}}+
	0.5*L_{\mathrm{Focal}},
	\]

	\subsection{Формирование бинарной маски}
	
	При необходимости карта логитов приводится к размеру входного изображения с помощью билинейной интерполяции.Бинарная маска формируется с порогом $\tau = 0.5$:
	\[
	\hat{y}_i =
	\begin{cases}
		1, & p_i \geq \tau,\\
		0, & p_i < \tau.
	\end{cases}
	\]
	
	Таким образом,изменение размера выполняется до преобразования логитов в вероятности и бинаризации.Дополнительная постобработка не применялась.Правила формирования итоговой маски задаются одинаково для всех сравниваемых моделей.
	
	% ============================================================
	\section{Заключение}
	
	После получения результатов
	
	% ============================================================
	\begin{thebibliography}{9}
		\bibitem{ronneberger2015} Ronneberger O., Fischer P., Brox T. U-Net: Convolutional Networks for Biomedical Image Segmentation. \textit{Medical Image Computing and Computer-Assisted Intervention --- MICCAI 2015}, 2015, vol.~9351, pp.~234--241. DOI: 10.1007/978-3-319-24574-4\_28.
		
		\bibitem{zhou2018} Zhou Z., Siddiquee M.~M.~R., Tajbakhsh N., Liang J. UNet++: A Nested U-Net Architecture for Medical Image Segmentation. \textit{Deep Learning in Medical Image Analysis and Multimodal Learning for Clinical Decision Support}, 2018, vol.~11045, pp.~3--11. DOI: 10.1007/978-3-030-00889-5\_1.
		
		\bibitem{xie2021} Xie E., Wang W., Yu Z., Anandkumar A., Alvarez J.~M., Luo P. SegFormer: Simple and Efficient Design for Semantic Segmentation with Transformers. \textit{Advances in Neural Information Processing Systems}, 2021, vol.~34, pp.~12077--12090. URL: https://arxiv.org/abs/2105.15203.
		
		\bibitem{wang2020} Wang J., Sun K., Cheng T., Jiang B., Deng C., Zhao Y., Liu D., Mu Y., Tan M., Wang X., Liu W., Xiao B. Deep High-Resolution Representation Learning for Visual Recognition. \textit{IEEE Transactions on Pattern Analysis and Machine Intelligence}, 2021, vol.~43, no.~10, pp.~3349--3364. DOI: 10.1109/TPAMI.2020.2983686.
		
		\bibitem{lee2024} Lee Y., Kim S.~B. Weakly Supervised Image Segmentation for Detecting Defects from Scanning Electron Microscopy Images in Semiconductor. \textit{IEEE Access}, 2024, vol.~12, pp.~184896--184908. DOI: 10.1109/ACCESS.2024.3513873.
				
		\bibitem{hasic2025} Hasić V., Kofler C., Krivić S. Towards Explaining SEM Defect Image Classification. \textit{ECAI 2025. Frontiers in Artificial Intelligence and Applications}, 2025, vol.~413, pp.~5463--5470. DOI: 10.3233/FAIA251487.
		
	\end{thebibliography}
	
\end{document}